% !LW recipe = latexmk (xelatex)
\documentclass[12pt,a4paper,UTF8]{ctexart}
\usepackage{SJTUReport}
\usepackage{booktabs} % 三线表
\usepackage{amsmath}
\usepackage{amssymb}

%%-------------------------------正文开始---------------------------%%
\begin{document}

% %-----------------------封面--------------------%%
\cover

%------------------摘要-------------%%
\newpage
\begin{abstract}

面向高速公路长途出行场景，本文研究沿线换电站网络中预约与随机两类异质用户的实时运营优化问题。长途电动汽车在一次行程中可能需要多次换电，每次换电在取走满电电池的同时退回一块低电量电池，使用户的换电站序列选择与沿线各站的电池库存形成时空耦合。为此，本文构建模型预测控制（MPC）与强化学习（RL）相结合的分层协同框架：MPC 在滚动预测域内联合优化预约用户剩余行程的换电站路径与请求的槽位匹配，RL actor 为各站每个槽位电池生成请求充电功率，critic 则通过分段线性状态价值函数，统一评估期末库存与未完成换电请求对后续运营净收益的影响。模型以连续时间刻画到站、等待、超时、充电完成与瞬时换电事件，允许仍在等待的请求带入下一轮，并按真实发生时刻对服务、违约与路径发布结算一次。

\end{abstract}
\newpage

\thispagestyle{empty} % 首页不显示页码

%%--------------------------目录页------------------------%%
\newpage
\tableofcontents

%%------------------------正文页从这里开始-------------------%
\newpage

%可选择这里也放一个标题
\begin{center}
   \title{ \Huge \textbf{{考虑异质用户的高速公路换电站运营 MPC-RL 分层协同优化}}}
\end{center}


\section{Introduction}
随着电动汽车在高速公路中长距离出行场景中的普及，沿线补能设施的运营能力逐渐成为影响用户出行可靠性与系统运行效率的关键因素。相比传统充电方式，换电模式可以缩短车辆补能时间。在高速公路场景下，长途出行的电动汽车往往需要沿行程多次换电：车辆沿道路单向行驶且不能折返，前一次换电的站点选择直接改变后续可达站点集合与剩余行程的换电站路径，各次换电决策因此相互关联。

与此同时，不同车辆到达时的电池荷电状态（SOC）各不相同：SOC 通过影响剩余续航里程，决定车辆首个可达换电站以及整条换电站序列；而每次换电在换出满电电池的同时退回一块低 SOC 电池，退回电池须经重新充电后才能再次服务。因此，一位用户的站点选择不仅影响自身后续行程，还会改变相关站点未来的电池库存状态，使用户路径与多站点电池库存产生时空耦合。高速公路换电站运营由此成为一个具有用户车辆电池状态和多站点电池状态时空耦合、动态演化特征的序贯决策问题。

实际运营中存在两类异质用户：预约用户和随机到达用户。预约用户在运营日前提交计划到达高速入口的时间、行程起点和终点以及到达 SOC，运营者据此完成接纳决策并发布日前换电路径；随机用户直接到站提出请求，其到达时刻与 SOC 只能预测。预约用户的实际到达信息可能偏离预约内容；车辆进入高速公路后，其未执行的剩余行程的换电站路径需要随实时位置、实时 SOC 和预计到达时间不断更新；到站时若暂无满电电池可用，用户在站内等待，超过其等待耐心后流失。考虑到分时电价使站内充电成本随时间变化，如何协调预约用户履约、随机用户接纳以及站内电池充电调度，是高速公路换电网络运营中的核心问题。

针对上述问题，本文构建 MPC--RL 分层协同优化框架。换电站的实时运营中存在两类性质不同但相互耦合的决策：一类是预约用户剩余换电路径的动态调整决策，即在有限预测域内根据未到用户的预约信息和在途用户的实时状态调整其未执行的换电站序列；另一类是各站每个槽位电池的充电功率决策。剩余路径调整具有离散组合结构，并同时受到车辆续航、路径连通性、预约履约和多站点库存约束；随机需求预测与实际到达又随时间不断更新。模型预测控制（Model Predictive Control, MPC）能够在有限预测域内联合评估这些约束和未来库存影响，并在每次获得新观测后重新优化而只执行当前外层区间，因而适合处理该类动态路径调整问题。充电功率的影响则具有明显的跨时段累积性：当前功率不仅决定当期充电成本，还会改变未来各站每个槽位电池的 SOC、满电电池供给能力和后续服务水平。本文采用强化学习（Reinforcement Learning, RL）学习每个槽位的充电功率控制策略，使其能够在长期交互过程中根据需求、电价和电池状态自适应调整充电行为。如\autoref{fig:MPC-RL}所示，RL actor 输出预测域内的每个槽位的请求充电功率并作为 MPC 的已知参数，MPC 在给定功率下联合优化剩余路径与请求的槽位匹配，critic 向 MPC 提供预测域之外的长期价值信息。

\begin{figure}[!htbp]
  \centering
  % FIGURE_PROMPT_BEGIN: fig:MPC-RL
  % 绘制一张适用于运筹优化与智能交通学术论文的 MPC–RL 协同框架图。白色背景，扁平化矢量风格，横向 16:9，配色仅使用深蓝、青绿色、橙色和中性灰，无渐变、阴影、3D 或装饰图标。
  % 图中从左到右展示：当前观测状态（每个槽位电池 SOC、预约与随机等待队列及剩余耐心、未到和在途预约信息、实时位置与 ETA、电价和随机需求预测）；基于已知信息预先确定的计划换电路径与计划执行轨迹；RL actor 输出预测域内每个槽位的请求功率，critic 提供固定参数的分段线性价值模型，在 MPC 内评价随决策变化的终端状态；MPC 在给定 RL 信息后优化预约用户剩余路径和站内连续事件安排；确定性执行器在当前外层区间内按真实到站、充电完成、换电和超时事件连续执行；输出实际 reward、更新后的电池状态、等待队列和已发布剩余路径，并反馈至下一轮。
  % 明确区分外层固定间隔决策与区间内部连续事件，强调计划换电路径先确定、随后 actor 滚动推演、最后 MPC 求解，不出现 RL–MPC 内循环。所有文字使用简洁中文，字体统一，箭头方向明确，模块不超过六个，适合缩放后在 A4 论文中阅读。输出优先为可编辑 SVG/PDF 或高分辨率透明背景图片。
  % FIGURE_PROMPT_END: fig:MPC-RL
  \fbox{\parbox[c][0.52\textwidth][c]{0.92\textwidth}{%
    \centering\small Figure placeholder}}
  \caption{MPC--RL 协同决策、连续事件执行与状态反馈框架}
  \label{fig:MPC-RL}
\end{figure}

进一步地，本文利用 RL critic 提供预测时域之外的长期价值信息，以缓解有限预测时域导致的末端效应。传统 MPC 通常通过设置固定终端库存下界避免时域末端过度消耗满电电池，但该方法难以同时反映库存 SOC 分布和未来换电请求的时间差异。本文以预测期末之后的运营净回报训练状态价值函数，统一评估期末库存与未完成请求的长期影响；未来请求保留其到达时间与等待期限，在后续实际服务时才消耗满电库存。为便于嵌入 MPC，critic 采用分段线性结构，其参数在每轮求解中固定，输入的终端状态则随本轮路径与服务决策变化。

因此，本文提出的 MPC--RL 框架基于决策属性进行分层：MPC 负责处理强约束、组合结构明显且需要可解释性的预约用户剩余路径动态调整问题；RL 负责学习具有长期影响的各槽位充电功率策略，并通过 critic 向 MPC 提供后续运营价值信息。二者通过每个槽位 SOC 状态、每个槽位的请求充电功率、在途预约履约事件和终端价值项相互连接，从而在保证预约用户服务质量要求和运营约束可行性的同时，使有限时域 MPC 能够近似感知预测时域之外的长期影响。

% CONTRIBUTIONS_PLACEHOLDER_BEGIN
% 待作者补充本文贡献
% CONTRIBUTIONS_PLACEHOLDER_END

本文其余部分安排如下：第 2 节给出问题描述；第 3 节建立滚动时域下的 MPC 模型，包括基于 SOC 分档的离线候选网络、日前基准路径生成和在线优化模型；第 4 节建立充电控制的马尔可夫决策过程并介绍 RL 算法；第 5 节给出数值实验；第 6 节总结全文。全文将站点规定的满电 SOC 阈值归一化为 $1$；因此，“满电”表示达到该运营目标值，并不要求等于电芯的电化学上限。






\section{Problem Statement}

考虑由多个 O-D 对及沿线换电站构成的单向高速公路换电网络。对于任一 O-D 对 $p$，车辆从入口 $o^p$ 进入高速公路后沿下游方向顺序行驶且不能折返，最终从出口 $d^p$ 离开；一次长途行程可能需要在沿线换电站集合 $\mathcal S^p$ 中的多个站点换电。全系统站点集合为 $\mathcal S=\bigcup_{p\in\mathcal P}\mathcal S^p$，不同 O-D 对的沿线站点可以重叠。站点 $i$ 配备槽位集合 $\mathcal B_i$，槽位中的电池以归一化 SOC 刻画电量状态。每次换电从站点取走一块 SOC 为 $1$ 的满电电池，并退回一块带有剩余电量的电池，退回电池入槽充电后可再次服务。因此，某位用户的站点选择不仅决定其后续可达站点，还会改变相关站点未来的电池库存，使用户路径与多站点库存产生时空耦合。

系统服务预约用户和随机到达用户。预约用户在运营日前提交预约信息，包括 O-D 对、预计进入高速公路的时刻及进入时的 SOC，运营者据此决定是否接纳预约，并为已接纳用户制定和发布日前换电路径。随机用户在到达换电站时提出请求，运营者只掌握其预测到达信息。按运行状态，已接纳预约用户分为两类：尚未进入高速公路的未到用户，以及已进入高速公路的在途用户；已经迟于预约入口时刻但尚未实际进入的用户仍属于未到集合，不能仅按原预约时刻将其删除。对每个 $p\in\mathcal P$，记 $\mathcal K_\ell^p=\mathcal K_\ell^{\mathrm{fut},p}\mathbin{\dot\cup}\mathcal K_\ell^{\mathrm{enr},p}$，其中 $\dot\cup$ 表示不交并；已完成行程或已失败退出的用户不再纳入。未到用户的有效入口 SOC 记为 $\widetilde{\mathrm{soc}}_{A,\ell}^{p,k}$，其初值为预约入口 SOC，并在获得新的入口前观测后更新；用户实际进入后，系统改用其实时位置、实时 SOC 和到达下游节点的预计到达时间。到站时若暂无满电电池可用，用户在站内等待：任一请求 $r$ 自其到站时刻 $t_r^{\mathrm{arr}}$ 起最长等待 $\Delta t$，开始服务时刻 $t_r^{\mathrm{sw}}$ 须满足 $t_r^{\mathrm{arr}}\le t_r^{\mathrm{sw}}\le t_r^{\mathrm{ddl}}$，其中 $t_r^{\mathrm{ddl}}=t_r^{\mathrm{arr}}+\Delta t$ 为等待截止时刻；截止时刻仍未获得服务的预约记为一次预约失败，随机请求则流失。换电视为瞬时事件。

系统以固定长度 $\sigma$ 的外层区间组织运营。以运营日统一时间原点，记
\[
t_q=q\sigma,
\qquad
\mathcal I_q=[t_q,t_{q+1}),
\qquad
\mathcal T_\ell=\{\ell,\ldots,\ell+H-1\}.
\]
在每个区间起点 $t_\ell$，系统以未来 $H$ 个区间为预测域，即 $\mathcal T_\ell$ 对应连续时间范围 $[t_\ell,t_{\ell+H})$。本轮预测采用固定的连续到站 ETA，不传播预测等待造成的下游延误，下一轮再依据实测更新，区间索引只用于请求归组与电价等分段常数数据。服务顺序遵循两条规则：预约请求优先于随机请求；同类请求内部按原到站时刻升序，同刻到达的同类请求采用预先给定的顺序。系统不为尚未到站的请求扣留当前可用电池；未服务完的等待请求带入下一轮，保留其原到站时刻与原截止时刻，不重新计时。

运营者需要协调预约用户剩余行程的换电站路径与各站每个槽位电池的充电功率，从而在履行已接纳预约的前提下，权衡换电收益、购电成本、路径调整成本和长期库存价值。本文采用分层 MPC--RL 框架：RL 根据需求、电价和库存状态生成预测域内的每个槽位的请求充电功率，并向 MPC 提供分段线性终端价值模型；MPC 在给定请求功率下联合优化剩余路径与请求的槽位匹配。每轮优化后，系统发布发生变化的在途用户剩余行程的换电站路径，只执行首个外层区间内的充电与服务安排，随后根据真实到站、充电完成、换电和超时事件更新状态并进入下一轮求解。

\autoref{tab:notation}列出正文反复使用的核心符号；仅局部使用的符号在首次出现处就近说明。RL 的状态、动作和算法符号在第 4 节定义。

\begin{table}[!htbp]
  \centering
  \caption{正文主要符号}
  \label{tab:notation}
  \small
  \setlength{\tabcolsep}{3pt}
  \renewcommand{\arraystretch}{0.98}
  \begin{tabular}{@{}lp{11cm}@{}}
    \toprule
    \textbf{符号} & \textbf{含义} \\
    \midrule
    \multicolumn{2}{l}{\textit{集合与索引}} \\
    $\sigma,t_q,\mathcal I_q,\mathcal T_\ell,\Delta t$ & 外层区间长度与起点、半开区间、含 $H$ 个区间的预测域索引集，以及最大等待时间 \\
    $\mathcal P,p;\ \mathcal S^p,\mathcal S;\ \mathcal B_i,b$ & O-D 对、沿线与全系统站点、库存槽位的集合和索引；槽位不是永久电池编号 \\
    $\mathcal K_\ell^{\mathrm{fut},p},\mathcal K_\ell^{\mathrm{enr},p},\mathcal K_\ell^p$ & 未到、在途预约用户集合及其并集；迟到但尚未实际进入者仍属于未到集合 \\
    $o^p,d^p,o_\ell^{\mathrm{cur},p,k},\widetilde{\mathrm{soc}}_{A,\ell}^{p,k},\mathrm{soc}_{A,\ell}^{p,k}$ & 入口、出口、实时虚拟起点、未到用户有效入口 SOC 与在途用户实时 SOC \\
    $\mathcal A^{p,k},t_{A,i}^{p,k},t_r^{\mathrm{arr}},t_r^{\mathrm{ddl}},t_r^{\mathrm{sw}}$ & 个体剩余候选弧集、连续 ETA，以及请求的到站、截止和开始服务时刻 \\
    \midrule
    \multicolumn{2}{l}{\textit{参数与状态}} \\
    $D_{i,j}^{p},v_{i,j}^{p},\underline D,\underline s_d^p$ & 节点间距离与 SOC 消耗、最小换电间距和出口最低 SOC \\
    $c^{ch}_{i,q},\pi_{i,q}^{\mathrm{sw}},\overline p_i,\overline P_i$ & 购电价、服务价格、单槽功率上限和站级总功率上限 \\
    $E_B,\eta_i,\widehat P_{i,b,q},T_{i,b,q}^{\mathrm{ch}},S_{i,b}(t)$ & 电池额定能量、充电效率、RL 请求功率、槽位在时段内的实际充电时长和槽位 SOC \\
    $\bar{\mathbf y}_{A}^{p,k},\mathbf y_{A}^{\mathrm{pub},p,k}$ & 日前初始发布路径与最近一次实际发布的剩余路径 \\
    $c^{adj},c^{\mathrm{fail}},\beta$ & 单次路径调整成本、单次预约失败成本与 RL 终端价值权重 \\
    \midrule
    \multicolumn{2}{l}{\textit{决策变量与主要辅助量}} \\
    $y_{i,j}^{p,k},x_{A,i}^{p,k},d_{p,k}$ & 剩余路径、站点访问量与路径差异指示量 \\
    $s_r,s_{r,q},f_r,\omega_r,a_r$ & 服务及其所属时段、失败或流失、期末等待和请求激活的辅助指示量 \\
    $z_{r,b};\ \mathcal C_i^A,\mathcal E_i,\mathcal E_i^A$ & 请求与槽位的匹配决策；全行程新预约事件、域内候选请求及其中预约请求的集合 \\
    $s_{\ell+H},\Phi^{\mathrm{RL}},w_h^{\mathrm{pend}}$ & 本轮决策对应的终端状态、后续运营价值及期末未完成换电指示量 \\
    \bottomrule
  \end{tabular}
\end{table}


\section{MPC Formulation}

设当前决策区间为 $\mathcal I_\ell$。在时刻 $t_\ell$，系统读取：各站每个槽位电池的观测 SOC；自上一轮遗留、仍未超时的站内等待请求；在途预约用户的实时位置、实时 SOC 与到达下游站点的预计到达时间；已发布且尚未执行的预约用户剩余行程的换电路径；以及未来 $H$ 个区间的随机需求预测与电价。

随后，系统按固定顺序构造本轮求解输入。首先确定计划换电路径：在途用户取最近一次实际发布路径；未到用户优先沿用上一轮保留的计划换电路径。其次，以RL actor 的策略，在计划换电路径下滚动生成预测域内的每个槽位的充电功率 $\widehat P_{i,b,q}$。

在此基础上，系统为每位预约用户构造剩余行程的候选路径网络，并由 MPC 更新剩余行程的换电路径安排。

求解后，系统立即发布发生变化的在途用户剩余行程的换电站路径并结算一次路径调整成本，新路径作为下一轮的调整基准；未到用户的优化路径暂不发布。系统只执行首个外层区间 $\mathcal I_\ell$：区间内按真实到站、充电完成、换电和超时事件连续推进；区间结束时，每个槽位 SOC、仍未超时的等待队列与已发布剩余路径进入下一轮。上述过程如\autoref{fig:rolling_horizon}所示。

\begin{figure}[!htbp]
  \centering
  % FIGURE_PROMPT_BEGIN: fig:rolling_horizon
  % 绘制一张适用于学术论文的横向滚动时域示意图。白色背景、简洁矢量风格，宽高比约 2:1，使用深蓝表示预测域、橙色表示当前执行区间、绿色表示带入下一轮的等待请求、灰色表示已执行历史。
  % 顶部为时间轴，清晰标出 $t_{\ell-1}$、$t_\ell$、$t_{\ell+1}$ 和 $t_{\ell+H}$，预测域采用半开区间 $[t_\ell,t_{\ell+H})$，突出只执行首个外层区间。下方设置三条泳道：第一条为未到预约用户，显示其内部路径可滚动优化但仍相对初始发布路径评价；第二条为在途预约用户，显示从实时虚拟起点重优化剩余行程的换电站路径，路径变化后发布并更新比较基准；第三条为站内等待队列，显示原始到站时刻和截止时刻跨滚动边界保持不变。
  % 当前执行区间内部使用少量时间标记展示连续到站、充电完成、立即换电和超时事件；预测终点处显示尚未截止的请求传入下一轮。不得出现“已发布路径固定不再优化”“整时段统一换电”或“随机用户必须当期完成”等旧逻辑。文字精炼、层级清楚、无大段说明，适合 A4 论文单栏宽度阅读。输出优先为可编辑 SVG/PDF 或高分辨率图片。
  % FIGURE_PROMPT_END: fig:rolling_horizon
  \fbox{\parbox[c][0.46\textwidth][c]{0.92\textwidth}{%
    \centering\small Figure placeholder}}
  \caption{滚动时域中的剩余路径优化、连续事件执行与等待请求的跨轮传递}
  \label{fig:rolling_horizon}
\end{figure}


\subsection{Offline Candidate networks by SOC Range}

在每个决策时刻，需要安排预测域内每位预约用户剩余行程的换电路径。如何高效、清晰地表示解的结构并设计决策变量是首要问题。本文借鉴加油站选址模型中 expanded network 的思路，用一条入口--出口路径表示该序列：对于 O-D 对 $p$，若路径为
$o^p\rightarrow i_1\rightarrow\cdots\rightarrow i_r\rightarrow d^p$，则用户依次在换电站 $i_1,\ldots,i_r$ 换电，一条弧跨过的沿线站点均不被访问。为避免在每个滚动时刻重复枚举全部下游组合，本文在每轮优化前先生成候选路径网络。不同的用户出发 SOC 会影响车辆的续航里程，从而影响可达的换电站集合。我们将用户出发 SOC 划分为若干个离线区间，并在每个区间上生成相应的候选网络。

对每个 O-D 对 $p$，节点集合为
$\mathcal V^p=\{o^p\}\cup\mathcal S^p\cup\{d^p\}$，所有节点均按高速公路下游方向排列。参数 $D_{i,j}^p$ 和 $v_{i,j}^p>0$ 分别表示从 $i$ 到下游节点 $j$ 的行驶距离和 SOC 消耗。

考虑到现实场景中为用户安排两个距离过近的连续站点换电并不合理，同时也为了减少网络规模，给定最小换电间距 $\underline D>0$，生成的网络在保证方案可行性的基础上，尽量排除从出发位置到首个换电站以及任意两个连续访问换电站之间距离小于 $\underline D$ 的方案。指向出口的末段行程不受该间距约束，但车辆到达出口时的 SOC 不得低于 $\underline s_d^p$。

对每个 SOC 区间 $h$ 和 O-D 对 $p$，离线网络按以下步骤生成。

\begin{enumerate}
  \item[\textbf{Step 1.}] \textbf{生成原始弧。}若车辆以该 SOC 区间的上限电量能够从入口 $o^p$ 到达下游换电站 $j$，则连接二者；若 $v_{i,j}^p\le1$，则连接两个依次位于下游的换电站 $i$ 和 $j$；若 $v_{i,d^p}^p+\underline s_d^p\le1$，则连接换电站 $i$ 与出口 $d^p$。

  \item[\textbf{Step 2.}] \textbf{生成完整可行路径。}在原始网络中枚举所有从 $o^p$ 到 $d^p$ 的完整路径。

  \item[\textbf{Step 3.}] \textbf{剪枝过近换电弧。}从上述完整路径中找出距离小于 $\underline D$ 的首段弧和站间弧，如果现有的完整路径中仍存在至少一条不包含该弧的方案，则删除该弧，否则保留。

  \item[\textbf{Step 4.}] \textbf{形成候选网络。}将最终保留的完整路径所包含的弧合并为 $\mathcal A^{h,p}$，候选网络由节点集 $\mathcal V^p$ 和弧集 $\mathcal A^{h,p}$ 组成。
\end{enumerate}

离线生成的候选网络刻画了站点之间的下游连接关系，对尚未进入高速的用户，直接按其有效入口 SOC $\widetilde{\mathrm{soc}}_{A,\ell}^{p,k}$ 所在的档位选取对应的离线网络。对已经在途的用户，将当前位置设置为虚拟入口，将当前 SOC 看作有效入口SOC。在线构造时，再按用户的实际有效 SOC 删除不可达的首段弧，并保证所有出口入弧满足出口最低 SOC；在途用户已执行的路径前缀不再参与优化。已在站内等待用户的当前换电作为固定事件保留，具体衔接方式见 Optimization Model 小节。针对每位预约用户选取的候选弧集记为 $\mathcal A^{p,k}$。


\subsection{Day-Ahead Baseline Path Generation}

用户提交预约信息
$\left(p,\bar t_A^{p,k},\overline{\mathrm{soc}}_A^{p,k}\right)$
后，运营者依据其行程、预计入口时刻和入口 SOC 检查能否提供服务，并为可接纳的请求生成日前基准路径。

预约请求按提交顺序处理。处理当前请求前，算法根据已接纳预约和日前随机需求预测，递推估计各站未来各时段的电池库存。日前充电功率由 RL 在日前预测状态上逐区间生成。电池未满时按当前区间的日前功率充电，达到 SOC $1$ 后立即停止充电，换电退回的低 SOC 电池入槽后立即恢复充电。

对 O-D 对 $p$ 的预约用户 $k$，基准路径按以下规则生成：
\begin{enumerate}
  \item 车辆从入口出发时的 SOC 为 $\overline{\mathrm{soc}}_A^{p,k}$，每次换电后恢复为 $1$。下一节点沿候选弧集选择，因此首段和站间行程均满足续航要求。
  \item 若车辆能够从当前位置直接到达出口，且到达 SOC 不低于 $\underline s_d^p$，则下一站选择出口；否则，仅保留预计到站时刻已经有满电电池可用的候选站。
  \item 候选站按预计满电电池余量从高到低排序，余量相同时优先选择更靠近出口的站点。选定站点后，在相应时段为该用户预留一块满电电池，并更新该站后续时段的预计库存；随后重复上述过程。
\end{enumerate}

若任一步不存在满足库存与下游可达性要求的候选站，即不能生成完整入口--出口路径，则拒绝该请求；只有成功生成完整路径的用户才进入已接纳集合。对已接纳用户，若基准路径经过弧 $(i,j)$，则令 $\bar y_{i,j}^{p,k}=1$，否则为 $0$。该路径随预约结果一并发布，既作为用户的初始发布路径，也作为其未到阶段路径调整的固定比较基准。


\subsection{Optimization Model}

\paragraph{集合与变量。}
在时刻 $t_\ell$，对每个 $p\in\mathcal P$ 和 $k\in\mathcal K_\ell^p$，记剩余候选网络为 $(\mathcal V^{p,k},\mathcal A^{p,k})$，起点 $o_\ell^{p,k}$ 对未到用户为入口，对在途用户为实时虚拟起点。候选弧已按本轮有效 SOC 筛选，保证首段可达、站间行程可由满电完成且到达出口时满足最低 SOC。若用户已在站内等待，其当前换电必须保留：虚拟起点仅通过一条零行驶距离的固定弧连接当前站，后续路径从该站换电成功后出发；该固定弧对应的请求直接沿用遗留请求，不另行生成。

记 $\mathcal C_i^A$ 为剩余候选弧 $(j,i)\in\mathcal A^{p,k}$ 对应的新预约事件集合，元素 $r=(p,k,j,i)$ 包含用户及入弧信息。它覆盖整个剩余行程，但不包括上述已由遗留请求表示的换电。域内部分为
\[
\mathcal C_{i,\ell}^A
=\{r\in\mathcal C_i^A:t_\ell\le t_r^{\mathrm{arr}}<t_{\ell+H}\}.
\]
令 $\mathcal E_{i,\ell}^{R,\mathrm{new}}$ 为域内新到达的预测随机请求集，$\mathcal W_{i,\ell}$ 为期初仍在等待的遗留请求集，$\mathcal W_{i,\ell}^A\subseteq\mathcal W_{i,\ell}$ 为其中的预约请求。于是
\begin{equation}
\mathcal E_i=\mathcal C_{i,\ell}^A\mathbin{\dot\cup}\mathcal E_{i,\ell}^{R,\mathrm{new}}\mathbin{\dot\cup}\mathcal W_{i,\ell},
\qquad
\mathcal E_i^A=\mathcal C_{i,\ell}^A\mathbin{\dot\cup}\mathcal W_{i,\ell}^A
\subseteq\mathcal E_i,
\label{eq:request_sets}
\end{equation}
其中 $\dot\cup$ 表示不交并。上述集合均在求解前固定；路径决定候选请求是否激活，而不改变索引集合。新预约事件使用本轮更新且不早于 $t_\ell$ 的 ETA。各请求的到站时刻 $t_r^{\mathrm{arr}}$、截止时刻 $t_r^{\mathrm{ddl}}=t_r^{\mathrm{arr}}+\Delta t$ 及退回 SOC $\rho_r\in[0,1)$ 均为已知参数，遗留请求保留原到站与截止时刻。

本轮有两类决策变量：
\begin{equation}
\begin{aligned}
y_{i,j}^{p,k}&\in\{0,1\},
&&p\in\mathcal P,\ k\in\mathcal K_\ell^p,\ (i,j)\in\mathcal A^{p,k},\\
z_{r,b}&\in\{0,1\},
&&i\in\mathcal S,\ r\in\mathcal E_i,\ b\in\mathcal B_i.
\end{aligned}
\label{eq:decision_variables}
\end{equation}
$y_{i,j}^{p,k}=1$ 表示剩余路径经过弧 $(i,j)$，是核心决策；$z_{r,b}=1$ 表示请求 $r$ 换取槽位 $b$ 中的电池，并将退回电池放入同一槽位。用户对满电槽位无偏好，但各槽位的给定功率可能不同，因而不同匹配可能影响后续运营；只有相关状态与参数可交换时，槽位置换才产生等价解。

请求激活量 $a_r$、服务量 $s_r$、失败或流失量 $f_r$、期末等待量 $\omega_r$ 及分时段服务量 $s_{r,q}$ 均为二元辅助量；服务时刻 $t_r^{\mathrm{sw}}$、槽位 SOC $S_{i,b}(t)$、充电时长 $T_{i,b,q}^{\mathrm{ch}}$、路径调整指示 $d_{p,k}$ 和终端状态 $s_{\ell+H}$ 由下述约束确定。RL 给出的功率 $\widehat P_{i,b,q}$ 与价值函数参数 $\theta$ 在本轮固定，不属于 MPC 决策变量。

\paragraph{目标函数。}
每次服务交付的净电量为 $E_B(1-\rho_r)$，按实际服务所属时段收费。预测域内的服务收益、充电成本、路径调整成本和预约失败成本分别为
\begin{align}
I&=E_B\sum_{i\in\mathcal S}\sum_{r\in\mathcal E_i}
\sum_{q\in\mathcal T_\ell}\pi_{i,q}^{\mathrm{sw}}(1-\rho_r)s_{r,q},
\label{eq:income}\\
C_{\mathrm{ch}}&=\sum_{i\in\mathcal S}\sum_{b\in\mathcal B_i}
\sum_{q\in\mathcal T_\ell}c^{ch}_{i,q}\widehat P_{i,b,q}T_{i,b,q}^{\mathrm{ch}},
\label{eq:chargingcost}\\
C_{\mathrm{adj}}&=c^{adj}\sum_{p\in\mathcal P}\sum_{k\in\mathcal K_\ell^p}d_{p,k},
\label{eq:adjustment_cost}\\
C_{\mathrm{fail}}&=c^{\mathrm{fail}}\sum_{i\in\mathcal S}\sum_{r\in\mathcal E_i^A}f_r.
\label{eq:reservation_failure_cost}
\end{align}
未到用户的路径调整相对日前发布路径作前瞻评价，在途用户则相对最近实际发布的剩余路径评价；真实调整费用只在发布时结算一次。

以第 4 节的固定 critic 评价本轮决策产生的终端状态：
\begin{equation}
\Phi^{\mathrm{RL}}=V_\theta(s_{\ell+H}).
\label{eq:rl_terminal_value}
\end{equation}
它估计期末之后的运营净回报，不重复计入域内已结算项目。给定权重 $\beta\ge0$，优化目标为
\begin{equation}
\max_{\mathbf y,\mathbf z}\quad
I-C_{\mathrm{ch}}-C_{\mathrm{adj}}-C_{\mathrm{fail}}+\beta\Phi^{\mathrm{RL}},
\label{eq:objective}
\end{equation}
其中所有辅助量均随路径和槽位匹配通过以下关系确定。

\paragraph{约束条件。}
（1）\textbf{路径与退回电量。}每位用户在其候选网络中选择一条从本轮起点到出口的路径：
\begin{equation}
\sum_{j:(n,j)\in\mathcal A^{p,k}}y_{n,j}^{p,k}
-\sum_{j:(j,n)\in\mathcal A^{p,k}}y_{j,n}^{p,k}
=\begin{cases}
1,&n=o_\ell^{p,k},\\
-1,&n=d^p,\\
0,&\text{其余节点},
\end{cases}
\label{eq:flow}
\end{equation}
对所有 $p\in\mathcal P$、$k\in\mathcal K_\ell^p$、$n\in\mathcal V^{p,k}$ 成立。候选网络无环，因此单位流确定一个访问站序；遗留预约用户通向当前等待站的固定弧取 $y=1$。

记 $s_{0,\ell}^{p,k}$ 为本轮起点的有效 SOC，未到用户取 $\widetilde{\mathrm{soc}}_{A,\ell}^{p,k}$，在途用户取实时 $\mathrm{soc}_{A,\ell}^{p,k}$。新预约事件 $r=(p,k,j,i)\in\mathcal C_i^A$ 的退回电量为
\begin{equation}
\rho_r=\rho_{A,j,i}^{p,k}
=\begin{cases}
s_{0,\ell}^{p,k}-v_{o_\ell^{p,k},i}^{p,k},&j=o_\ell^{p,k},\\
1-v_{j,i}^{p},&j\in\mathcal S^p,
\end{cases}
\label{eq:return_soc}
\end{equation}
其中首段耗电由本轮起点位置确定；随机与遗留请求的退回 SOC 分别取预测值与观测值。

（2）\textbf{路径调整。}对所有 $p\in\mathcal P$、$k\in\mathcal K_\ell^p$ 和 $i\in\mathcal S^p$，定义站点访问量与调整指示为
\begin{equation}
x_{A,i}^{p,k}=\sum_{j:(j,i)\in\mathcal A^{p,k}}y_{j,i}^{p,k},
\qquad
d_{p,k}=\max_{i\in\mathcal S^p}
\left|x_{A,i}^{p,k}-x_{A,i\mid\ell}^{\mathrm{ref},p,k}\right|,
\label{eq:path_adjustment_indicator}
\end{equation}
其中基准 $x_{A,i\mid\ell}^{\mathrm{ref},p,k}$ 对未到用户取日前路径访问量，对在途用户取最近发布路径的未执行部分。无入弧的站点访问量为 $0$；若沿线无换电站，则 $d_{p,k}=0$。单向路径的站序由所访站点确定，故增加或删除任一站点恰好记一次调整。

（3）\textbf{请求激活与跨站衔接。}对随机及遗留请求，$a_r=1$。对新预约事件 $r=(p,k,j,i)\in\mathcal C_{i,\ell}^A$，令 $\mathcal U_r$ 为该用户在上游站 $j$ 的候选预约换电集合，包括对应的遗留请求；首段事件无前序请求。对整段行程的候选换电 $h$，以 $s_h^{\mathrm{in}}$ 表示域内服务量：$h\in\bigcup_i\mathcal E_i$ 时取 $s_h$，其余取 $0$。于是
\begin{equation}
a_r=y_{j,i}^{p,k}g_r,
\qquad
g_r=\begin{cases}
1,&j=o_\ell^{p,k},\\
\displaystyle\sum_{h\in\mathcal U_r}s_h^{\mathrm{in}},&j\ne o_\ell^{p,k}.
\end{cases}
\label{eq:reservation_alive_chain}
\end{equation}
其中 $\mathcal U_r=\{h\in\mathcal C_j^A\mathbin{\dot\cup}\mathcal W_{j,\ell}^A:h\text{ 属于用户 }(p,k)\}$ 对站间事件成立，首段取 $\mathcal U_r=\varnothing$。流平衡保证被选中的前序换电至多一个。首站由路径激活，后续站须在上游成功换电后才有效；上游失败会使后续事件失效。记
\[
\mathcal E_{p,k}^A=\{r\in\textstyle\bigcup_i\mathcal E_i^A:r\text{ 属于用户 }(p,k)\},
\qquad
F_{p,k}=\sum_{r\in\mathcal E_{p,k}^A}f_r\in\{0,1\}.
\]
其中 $F_{p,k}$ 是该用户域内失败指示，保证一次失败只结算一次。为保持固定 ETA 下的时间顺序，对每个新预约事件 $r=(p,k,j,i)\in\mathcal C_{i,\ell}^A$ 还要求
\begin{equation}
\begin{aligned}
y_{j,i}^{p,k}(1-F_{p,k})&\le g_r,\\
a_r=s_h^{\mathrm{in}}=1&\ \Longrightarrow\
t_h^{\mathrm{sw}}\le t_r^{\mathrm{arr}},\qquad h\in\mathcal U_r.
\end{aligned}
\label{eq:arrival_causality}
\end{equation}
即未失败用户计划在域内到达的换电不能因前序尚未服务而被省略，且上游服务须先于下游到站。本轮固定 ETA 不传播等待延误，可能排除实际可顺延到站的方案；下一轮依据实测更新 ETA。

（4）\textbf{服务规则与槽位匹配。}对每个 $i\in\mathcal S$、$r\in\mathcal E_i$，服务与结局满足
\begin{equation}
\begin{aligned}
\sum_{b\in\mathcal B_i}z_{r,b}&=s_r,
&s_r+f_r+\omega_r&=a_r,\\
s_r=1&\ \Longrightarrow\
t_r^{\mathrm{sw}}\in[t_r^{\mathrm{arr}},t_r^{\mathrm{ddl}}]
\cap[t_\ell,t_{\ell+H}),\\
s_{r,q}&=\mathbf1\{s_r=1\text{ 且 }t_r^{\mathrm{sw}}\in\mathcal I_q\},
&s_r&=\sum_{q\in\mathcal T_\ell}s_{r,q}.
\end{aligned}
\label{eq:service_window}
\end{equation}
未服务请求不定义服务时刻，各 $s_{r,q}$ 均取 $0$。事件时刻包括请求到站、换电、等待截止、功率时段边界和由 SOC 轨迹递推得到的充满时刻，并包含预测域端点。在每个事件时刻，将已经激活且到站、尚未服务且未超过截止时间的请求组成队列，将 SOC 为 $1$ 的槽位组成可用集合。只要两者均非空，就立即按预约优先、类内先到先服务的顺序取队首，并由 $z$ 为其匹配一个当前满电槽位；同刻到达的同类请求采用预先给定的顺序。更新槽位与队列后重复这一过程，直至队列或满电槽位集合为空。因此不能保留可用电池等待未来请求，也不能为更高价格而推迟当前服务。

对同刻事件，先计入截至该时刻的充电和到站，再完成可行服务，最后将仍未服务且恰好截止的请求判为失败或流失。期末 $t_{\ell+H}$ 的事件留到下一轮处理。由此
\begin{equation}
f_r=(a_r-s_r)\mathbf1\{t_r^{\mathrm{ddl}}<t_{\ell+H}\},
\qquad
\omega_r=(a_r-s_r)\mathbf1\{t_r^{\mathrm{ddl}}\ge t_{\ell+H}\}.
\label{eq:outcome_boundary}
\end{equation}
这些事件规则与等式共同定义服务时刻和结局，它们不是可独立选择的接纳或排程决策。

（5）\textbf{SOC 与充电时长。}取 $S_{i,b}(t_\ell^-)=S_{i,b}^{\mathrm{obs}}$。若相邻事件时刻 $t_1<t_2$ 之间属于时段 $q$ 且没有换电，则
\begin{equation}
S_{i,b}(t_2^-)=\min\left\{1,\,
S_{i,b}(t_1^+)+\frac{\eta_i\widehat P_{i,b,q}}{E_B}(t_2-t_1)\right\}.
\label{eq:event_charging_transition}
\end{equation}
时间以小时计，功率与电池能量分别使用 kW 与 kWh。换电时只允许交付满电电池，被选槽位立即换入退回电池：
\begin{equation}
\begin{aligned}
z_{r,b}=1&\ \Longrightarrow\
S_{i,b}((t_r^{\mathrm{sw}})^-)=1,\quad
S_{i,b}((t_r^{\mathrm{sw}})^+)=\rho_r,\\
\sum_{r\in\mathcal E_i:\,s_r=1,\ t_r^{\mathrm{sw}}=t}z_{r,b}&\le1,
\qquad t\in[t_\ell,t_{\ell+H}).
\end{aligned}
\label{eq:slot_matching}
\end{equation}
上述关系对所有 $i\in\mathcal S$、$b\in\mathcal B_i$ 及相应请求成立，未被选中的槽位在该时刻无 SOC 跳变。同一槽位可以反复服务，但每次服务后必须重新充满。实际充电时长由 SOC 轨迹确定：
\begin{equation}
T_{i,b,q}^{\mathrm{ch}}=\int_{t_q}^{t_{q+1}}
\mathbf1\{S_{i,b}(t)<1,\ \widehat P_{i,b,q}>0\}\,\mathrm dt,
\qquad i\in\mathcal S,\ b\in\mathcal B_i,\ q\in\mathcal T_\ell.
\label{eq:charging_duration}
\end{equation}
RL 输出在求解前满足
\begin{equation}
0\le\widehat P_{i,b,q}\le\overline p_i,
\qquad\sum_{b\in\mathcal B_i}\widehat P_{i,b,q}\le\overline P_i,
\label{eq:power_projection}
\end{equation}
其中 $i\in\mathcal S$、$b\in\mathcal B_i$、$q\in\mathcal T_\ell$。这是给定功率的输入可行性条件；满电停充及换入后恢复充电改变实际用电量，而不重新优化功率。

（6）\textbf{终端状态。}令
\[
\overline{\mathcal D}_\ell^A
=\bigcup_{i\in\mathcal S}(\mathcal C_i^A\mathbin{\dot\cup}\mathcal W_{i,\ell}^A)
\]
为整段剩余行程的候选预约换电集合。对 $h\in\overline{\mathcal D}_\ell^A$，令 $\delta_h$ 为其路径选择指示：新事件取对应入弧的 $y$，遗留请求取 $1$。期末未完成指示为
\begin{equation}
w_h^{\mathrm{pend}}=\delta_h(1-F_{p(h),k(h)})(1-s_h^{\mathrm{in}}),
\qquad h\in\overline{\mathcal D}_\ell^A,
\label{eq:pending_delivery_alive}
\end{equation}
其中 $(p(h),k(h))$ 为请求所属用户。它表示路径选中、用户未失败且域内尚未完成的换电。终端状态 $s_{\ell+H}$ 由期末 SOC $\{S_{i,b}(t_{\ell+H}^-)\}_{i,b}$、$\omega_r=1$ 的等待请求、$w_h^{\mathrm{pend}}=1$ 的未完成预约换电及已知外生信息构造，并使用与 RL 训练一致的固定编码。已在等待的预约请求同时满足 $w_h^{\mathrm{pend}}=\omega_r=1$，只编码一次；其他未完成预约换电保留名义到达时间、截止时间、退回 SOC 和前序依赖，不视为已经到站的等待请求。由此，路径与槽位匹配通过服务过程和库存轨迹共同影响式\eqref{eq:rl_terminal_value}中的终端价值。


\section{RL Formulation}

RL 层用于学习每个槽位的充电功率控制策略，并通过分段线性 critic 为 MPC 提供期末之后的运营价值。本节将充电功率控制问题建模为马尔可夫决策过程，并给出本文采用的 RL 算法。

\subsection{Markov decision process}

\paragraph{State.}
在滚动决策时刻 $t_\ell$，RL agent 观测系统状态 $s_\ell$。除各站每个槽位电池的 SOC、电价和时间特征外，状态包含：未到与在途预约用户信息（在途用户含实时位置、实时 SOC 和最近发布的剩余行程的换电站路径）、站内预约与随机等待队列（含原到站时刻与剩余耐心）、最近一个外层区间的真实随机到达记录，以及未来逐请求随机需求预测。令
$\mathcal K_\ell^{\mathrm{fut},p}$ 表示时刻 $t_\ell$ 尚未进入高速公路的已接纳预约用户集合（覆盖范围超出当前预测域），用 $\bar{\mathbf y}_A^{p,k}$ 和 $\mathbf y_A^{\mathrm{pub},p,k}$ 分别汇集日前初始发布路径与最近实际发布路径的弧指示量，则预约信息集合为
\begin{equation}
\begin{aligned}
\mathcal Z_{A}
={}&
\left\{
\left(p,k,\bar t_A^{p,k},
\overline{\mathrm{soc}}_A^{p,k},
\bar{\mathbf y}_A^{p,k}\right):
k\in\mathcal K_\ell^{\mathrm{fut},p}
\right\}\\
&\cup
\left\{
\left(p,k,o_\ell^{\mathrm{cur},p,k},\mathrm{soc}_{A,\ell}^{p,k},
\mathbf y_A^{\mathrm{pub},p,k}\right):
k\in\mathcal K_\ell^{\mathrm{enr},p}
\right\}.
\end{aligned}
\label{eq:rl_reservation_information}
\end{equation}
站内等待队列记为
$\mathcal W_\ell=\{(r,i,t_r^{\mathrm{arr}},t_r^{\mathrm{ddl}})\}$，
其中元素为仍未服务且未超时的预约与随机请求。由于 $\mathcal Z_{A}$ 与 $\mathcal W_\ell$ 的元素数量随滚动时刻变化，分别使用 $f_A$、$f_W$ 将其编码为固定维度特征。再令 $f_B$、$f_R^{\mathrm{act}}$ 和 $f_R^{\mathrm{pred}}$ 分别编码每个槽位 SOC、真实随机到达记录和预测随机请求，则状态表示为
\begin{equation}
\begin{aligned}
s_\ell=\Big(&
f_B\!\left(\{S_{i,b}^{\mathrm{obs}}(t_\ell)\}_{i\in\mathcal S,b\in\mathcal B_i}\right),
f_A(\mathcal Z_{A}),
f_W(\mathcal W_\ell),\\
&
f_R^{\mathrm{act}}\!\left(
\{(i,r,t_R^{i,r},\mathrm{soc}_R^{i,r},z_{R,r,i}^{\mathrm{act}}):
r\in\mathcal R_{i,\ell-1}^{\mathrm{act}}\}_{i\in\mathcal S}
\right),\\
&
f_R^{\mathrm{pred}}\!\left(
\{(i,\nu,t_{R}^{i,\nu},\mathrm{soc}_{R}^{i,\nu}):
\nu\in\mathcal R_{i,q}\}_{i\in\mathcal S,q\in\mathcal T_\ell}
\right),
\{c^{ch}_{i,q}\}_{i\in\mathcal S,q\in\mathcal T_\ell},
\xi_\ell
\Big).
\end{aligned}
\label{eq:rl_state}
\end{equation}
其中 $\xi_\ell$ 表示时段、节假日或交通高峰等外生特征；$\mathcal R_{i,\ell-1}^{\mathrm{act}}$ 为上一区间真实到站的随机请求记录，$z_{R,r,i}^{\mathrm{act}}$ 为其实际服务结果。$f_R^{\mathrm{act}}$ 使用真实的到达时刻、SOC 和服务结果，$f_R^{\mathrm{pred}}$ 仅使用时刻 $t_\ell$ 可得的预测量，两者不共享数据字段。由于同一站点的槽位标签没有物理意义，$f_B$ 对各槽位 SOC 使用共享的分段线性基函数并在站内求和，以刻画 SOC 分布；actor 的每个槽位的输出仍采用相应的置换等变解码，参考推演按预先给定的满电槽位选择规则分配请求，MPC 则通过 $z_{r,b}$ 优化本轮槽位匹配。各状态编码采用仿射变换、ReLU 与求和构成，使 critic 从原始状态到价值输出保持分段线性。请求编码保留相对于当前状态时刻的到达时间与截止时间、退回 SOC、用户类别及剩余路径关系；同一次换电请求在预约信息与等待队列中只计一次。在 MPC 中，候选请求的固定属性可预先编码，再由 $w_h^{\mathrm{pend}}$、$\omega_r$ 及路径变量控制其是否计入；随决策变化的库存和名义服务记录仍需同步编码。训练与终端评价采用相同的特征定义和时间基准，后续需求预测覆盖期末之后，不能把所有未完成请求合并为无时间信息的总量。

\paragraph{Action.}
RL 的动作是各站点、各槽位在当前外层区间的连续请求充电功率：
\begin{equation}
a_\ell=
\{\widehat P_{i,b,\ell}\}_{i\in\mathcal S,b\in\mathcal B_i}.
\label{eq:rl_action}
\end{equation}
actor 先生成每个槽位的无约束请求，再将其欧氏投影到式\eqref{eq:power_projection}所定义的闭凸集上，并以该唯一投影作为动作 $\widehat P$；该确定映射同时满足单槽功率上限和站级总功率上限。策略概率密度及 PPO 比率均针对这一固定投影参数化所诱导的请求动作分布计算。

\paragraph{Rollout.}
为向 MPC 提供未来 $H$ 个区间的请求功率，系统从
$s_{\ell}^{\mathrm{ref}}=s_\ell$ 出发，利用同一 actor 的均值动作和推演状态递推进行滚动推演：
\begin{equation}
\{\widehat P_{i,b,q}\}_{i,b}
=
\mu_\phi(s_{q}^{\mathrm{ref}}),
\quad q\in\mathcal T_\ell,
\label{eq:rl_power_rollout}
\end{equation}
其中 $\mu_\phi$ 表示策略分布的均值。推演遵循固定的调用顺序：先按第 3 节规则确定计划换电路径，再进行 actor 滚动推演，最后将固定功率序列与固定参数的 critic 一并交给 MPC。计划换电路径的确定只使用求解前已知的位置、SOC、预计到达时间和候选网络；推演采用与实际执行相同的预约优先、等待、超时、充电与电池更换规则，并允许预约软失败。因此，推演状态 $s_q^{\mathrm{ref}}$ 与本轮待求的路径和服务决策无关，请求序列生成与本轮路径优化之间不存在循环依赖。构造后续推演状态时，区间 $q$ 的预测请求及其名义服务结果进入独立的名义历史编码器 $f_R^{\mathrm{nom}}$；该编码器与 $f_R^{\mathrm{act}}$ 输出维度相同，但不共享输入字段或语义，且只用于推演状态和预测终端状态，不写回真实状态。$f_R^{\mathrm{pred}}$ 所需的后续预测窗口由同一需求预测器的条件均值补齐，超出当前 $H$ 步输入的部分采用固定终端预测先验。整条请求功率序列作为本轮 MPC 参数；只有首项请求 $\widehat P_{i,b,\ell}$ 进入真实执行，下一滚动时刻重新生成序列。

\paragraph{Transition.}
给定 $s_\ell$ 和请求功率序列后，系统求解 MPC，得到逐用户最优路径
$\{y_{i,j}^{p,k*}\}$ 和槽位匹配计划 $\mathbf z^*=\{z_{r,b}^*\}$。在途用户中剩余行程的换电站路径发生变化者在本轮边界立即发布并结算一次路径调整成本，新路径作为下一轮推演与调整的基准；未到用户的路径只用于当前预测，不对用户发布。为避免内部预测路径进入真实执行，定义边界后的执行路径：在途用户使用本轮已经发布的最优剩余行程的换电站路径，未到用户仍使用日前发布路径 $\bar{\mathbf y}_{A}^{p,k}$；若后者在 $\mathcal I_\ell$ 内实际进入高速公路，仍先按该已发布路径执行，下一边界再作为在途用户更新。记两类用户在 $\mathcal I_\ell$ 内的真实到站记录分别为 $\mathcal H_{A,\ell}^{\mathrm{act}}$ 和 $\mathcal H_{R,\ell}^{\mathrm{act}}$，每条记录携带请求标识、站点、真实到站时刻和真实退回 SOC。

环境随后只在首个外层区间 $\mathcal I_\ell$ 内执行。确定性连续事件算子 $\operatorname{Exec}_\ell$ 的输入仅为真实到站记录、遗留等待队列、边界后已发布的执行路径、当前请求功率、观测 SOC 和槽位匹配计划 $\mathbf z^*$；其内部按充电完成、登记到站、队列分配、超时判定的同时刻顺序连续推进：
\begin{equation}
\begin{aligned}
(\mathbf S^{\mathrm{obs}}(t_{\ell+1}),\mathcal W_{\ell+1},
\mathbf s_{A,\ell}^{\mathrm{act}},\mathbf z_{R,\ell}^{\mathrm{act}},
\mathbf f_{A,\ell}^{\mathrm{act}})
=\operatorname{Exec}_\ell\!\Big(&
\mathbf S^{\mathrm{obs}}(t_\ell),\mathcal W_\ell,
\{\widehat P_{i,b,\ell}\}_{i,b},\\
&\{\mathbf y_{A,\ell}^{\mathrm{exec},p,k}\}_{p,k},\ \mathbf z^*,
\mathcal H_{A,\ell}^{\mathrm{act}},\mathcal H_{R,\ell}^{\mathrm{act}}
\Big).
\end{aligned}
\label{eq:actual_execution_operator}
\end{equation}
其中 $\mathbf y_{A,\ell}^{\mathrm{exec},p,k}=\mathbf y_A^{p,k*}$ 对在途用户成立，而对未到用户固定为 $\bar{\mathbf y}_{A}^{p,k}$。$\mathbf z^*$ 给出本轮 MPC 的匹配计划：真实队首请求有对应的预测匹配且计划槽位仍满电时，沿用该匹配；否则按预先给定的规则从当前满电槽位中选择，始终立即服务队首。实际执行事件只来自两类真实到站记录和遗留等待队列；预测集合不作为真实到达输入执行算子，预测偏差不触发停止执行，也不要求预测与实际终态强制相等。状态转移可表示为
\begin{equation}
s_{\ell+1}
\sim
\mathbb P\!\left(
\cdot\mid
s_\ell,a_\ell,
\mathrm{MPC}(s_\ell,\widehat{\mathbf P})
\right),
\label{eq:rl_transition}
\end{equation}
其中 $\mathbb P$ 表示随机到达与下一轮预测更新共同诱导的环境状态转移核，$\widehat{\mathbf P}=\{\widehat P_{i,b,q}:i\in\mathcal S,b\in\mathcal B_i,q\in\mathcal T_\ell\}$。这一定义与 MPC 的“优化完整预测域、只执行当前区间”保持一致。

\paragraph{Reward.}
RL reward 使用当前实际执行区间的已实现运营收益，并计入本轮实际发布路径产生的调整成本：
\begin{equation}
r_\ell
=
I_\ell^{\mathrm{act}}
-C_{\mathrm{ch},\ell}
-C_{\mathrm{adj},\ell}
-C_{\mathrm{fail},\ell}.
\label{eq:rl_reward}
\end{equation}
其中 $I_\ell^{\mathrm{act}}$、$C_{\mathrm{ch},\ell}$、$C_{\mathrm{adj},\ell}$、$C_{\mathrm{fail},\ell}$ 分别为区间 $\mathcal I_\ell$ 内真实发生的服务收益、充电成本、路径调整成本和预约超时成本，各项均按真实发生时刻结算，且跨轮只结算一次：服务收益按实际服务时刻所在时段的价格结算，充电成本按区间内的实际充电量结算，路径调整成本只在真实向在途用户发布新路径时结算，预约失败成本按截止时刻所在的区间结算；上一轮到站、本轮才服务的请求只在服务发生的区间计一次收入，到预测期末仍在等待的请求不作最终统计。MPC 目标中的预测随机收益、未来预约超时风险和未来区间成本不进入真实 reward；未到用户的前瞻调整成本也不在当期 reward 中重复实现。该 reward 使 actor 学习每个槽位的充电功率对当前成本、预约可靠性、路径稳定性和长期电池状态价值的共同影响。


\subsection{RL algorithm}

由于每个槽位的请求功率是连续动作，且需要在不确定需求环境下保持训练稳定性，本文采用 Proximal Policy Optimization (PPO) 作为 RL 层的训练算法。PPO 属于 actor-critic 方法，actor $\pi_{\phi}(a_t\mid s_t)$ 输出连续请求功率动作的随机策略，critic $V_{\theta}(s_t)$ 估计当前状态的长期价值。critic 采用式\eqref{eq:pwl_critic}的分段线性结构，直接学习 state value function，可由式\eqref{eq:rl_terminal_value}统一评价终端库存与未完成请求的后续影响；该价值对应 actor 与 MPC 组成的后续闭环控制下的期望折扣运营净回报，不另设逐槽位价值或未完成请求的独立扣减项。

$V_\theta$ 采用含 $M$ 个 ReLU 单元的分段线性结构：
\begin{equation}
V_\theta(s)
=\boldsymbol a^\top s+b
+\sum_{m=1}^{M}c_m
\max\{0,\boldsymbol w_m^\top s+d_m\}.
\label{eq:pwl_critic}
\end{equation}
价值网络及状态编码参数在每轮 MPC 求解中固定，仅其输入终端状态随本轮决策变化。

PPO 的 critic 通过最小化价值函数误差进行训练：
\begin{equation}
L_V(\theta)=
\mathbb{E}_t\left[
\left(V_{\theta}(s_t)-\hat{G}_t\right)^2
\right],
\label{eq:value_loss}
\end{equation}
其中 $\hat{G}_t$ 为由采样轨迹中后续实际运营 reward 计算的折扣回报目标，服务收入、充电费用及超时损失均按实际发生时刻计入。采样批次内固定供 MPC 使用的网络参数，批次结束后更新模型；训练样本需覆盖不同库存分布、请求到达时间与等待压力。actor 采用 clipped surrogate objective 更新策略。定义重要性采样比率
\begin{equation}
w_t^{\mathrm{IS}}(\phi)=
\frac{\pi_{\phi}(a_t\mid s_t)}{\pi_{\phi_{old}}(a_t\mid s_t)},
\label{eq:ppo_ratio}
\end{equation}
则策略优化目标为
\begin{equation}
L_{\mathrm{CLIP}}(\phi)=
\mathbb{E}_t\left[
\min\left(
w_t^{\mathrm{IS}}(\phi)\hat{A}_t,
\mathrm{clip}(w_t^{\mathrm{IS}}(\phi),1-\epsilon,1+\epsilon)\hat{A}_t
\right)
\right],
\label{eq:ppo_clip}
\end{equation}
其中 $\hat{A}_t$ 为优势函数估计，$\epsilon$ 为 clip 参数。优势函数由广义优势估计（Generalized Advantage Estimation, GAE）计算：
\begin{equation}
\hat{A}_t=\sum_{n=0}^{\infty}
(\gamma_{\mathrm{RL}}\lambda_{\mathrm{GAE}})^n
e_{t+n}^{\mathrm{TD}},
\quad
e_t^{\mathrm{TD}}=r_t
+\gamma_{\mathrm{RL}}V_{\theta}(s_{t+1})-V_{\theta}(s_t),
\label{eq:gae}
\end{equation}
其中 $\gamma_{\mathrm{RL}}$ 为折扣因子，
$\lambda_{\mathrm{GAE}}$ 为 GAE 参数。外层 RL 决策保持固定步长 $\sigma$，与区间内部的连续物理事件并存。

\paragraph{终端价值的嵌入。}在一次 MPC 求解中固定价值网络与编码参数，式\eqref{eq:pwl_critic}中的 ReLU 最大值可逐单元线性化。令 $\zeta_m=\boldsymbol w_m^\top s_{\ell+H}+d_m$，以 $u_m$ 表示其 ReLU 输出。给定有效界 $L_m\le\zeta_m\le U_m$，当 $L_m<0<U_m$ 时，引入 $\delta_m^V\in\{0,1\}$，有
\begin{subequations}
\label{eq:pwl_critic_embedding}
\begin{align}
u_m&\ge0,& u_m&\ge\zeta_m,\\
u_m&\le U_m\delta_m^V,&
u_m&\le\zeta_m-L_m(1-\delta_m^V),
\qquad m=1,\ldots,M.
\end{align}
\end{subequations}
其中 $L_m,U_m$ 由本轮可行状态的有效上下界推得；若 $U_m\le0$，直接取 $u_m=0$，若 $L_m\ge0$，直接取 $u_m=\zeta_m$。将式\eqref{eq:pwl_critic}中的最大值替换为 $u_m$，并对状态特征编码中的 ReLU 作同样处理，终端价值模块即整体表示为混合整数线性约束，其规模随编码及网络中需判定激活状态的单元数增加。

训练完成后，actor 在每个滚动时刻输出首期每个槽位的请求功率，并通过式\eqref{eq:rl_power_rollout}生成供 MPC 使用的预测域请求功率序列；实际仅执行首期请求功率。critic 的分段线性参数在本轮求解中固定，MPC 根据自身的路径与槽位匹配决策构造 $s_{\ell+H}$，通过式\eqref{eq:pwl_critic_embedding}联合确定终端价值。由此，RL 层同时提供请求功率轨迹和后续运营价值模型，MPC 则优化剩余路径与请求的槽位匹配，形成闭环协同机制；参考推演仅用于生成功率等求解输入，不预先固定本轮终端价值。


\section{Numerical experiments}

% NUMERICAL_EXPERIMENTS_PLACEHOLDER_BEGIN
% 待作者提供实验设置、数据与结果
% NUMERICAL_EXPERIMENTS_PLACEHOLDER_END


\section{Conclusion}

本文针对高速公路换电站网络中预约与随机两类异质用户的运营优化问题，构建了 MPC--RL 分层协同框架。在 MPC 层，本文以滚动时域组织运营：未到预约用户从入口出发、相对日前初始发布路径评价调整，在途预约用户从实时虚拟起点出发、相对最近发布的未执行站序评价调整，站内仍在等待的请求带入下一轮且保留原到站时刻与截止时刻；候选网络按 SOC 分档离线生成、在线裁剪，路径调整成本按实际发布结算一次。站内服务以连续时间刻画：请求在到达后最长等待 $\Delta t$，预约优先、类内按到站顺序服务，槽位电池严格在达到 SOC $1$ 后交付，充电在满电时停止、换入退回电池后立即恢复。在 RL 层，actor 输出每个槽位的请求充电功率作为 MPC 的已知参数，critic 通过分段线性状态价值函数，统一衡量期末库存与未完成请求对后续运营净回报的影响；计划换电路径先于 actor 滚动推演确定，随后在固定功率序列和网络参数下由 MPC 优化路径与槽位匹配，并据此计算服务结果与终端价值，避免请求序列与本轮路径决策之间的循环依赖。

本文的模型边界包括：车辆到达下游节点的预计到达时间为外生给定，未建立随机交通模型；换电视为瞬时事件，未区分服务开始与作业结束；终端价值采用有限维状态编码与分段线性 critic 近似，不能保证完整保留库存分布和多请求依赖，其误差及新增整数变量对求解效率的影响需通过实验评估。后续工作可围绕更精细的交通流建模、换电作业时长以及终端价值的精确刻画展开，并通过系统的数值实验评估框架在不同需求强度、电价模式和库存配置下的表现。


%%----------- 参考文献 -------------------%%
%在reference.bib文件中填写参考文献，此处自动生成

% 当前正文尚无 \cite 命令，添加正式引用后再启用下行参考文献列表。
% \reference

\end{document}
